\chapter{Configuration}
\label{ch:config}

Secrets live only in a \code{.env} file (permissions \code{0600}, never committed, never
logged). Non-secret settings live in a TOML config file the UI can re-read at runtime.

\section{smart-me authentication}
\label{sec:smartme-auth}

smart-me offers several authentication mechanisms. \prog uses the one designed for
headless devices.

\subsection{The available mechanisms}
\begin{itemize}
  \item \textbf{OAuth “Device (Client Credentials)” app} — described by smart-me as
        \emph{“a replacement for basic authentication in embedded devices\dots\ supports
        only the OAuth flow client credentials.”} This is a non-interactive
        machine-to-machine flow — exactly what a 24/7 daemon needs.
  \item \textbf{OAuth Confidential / Public apps} — use interactive flows
        (authorization-code~+~PKCE, device-code, implicit) that require a human to log in
        and consent. Not suitable for an unattended bridge.
  \item \textbf{HTTP Basic auth} — the smart-me account username and password.
\end{itemize}

\subsection{What \prog uses}
\begin{description}[style=nextline]
  \item[Primary — OAuth2 Client Credentials] The bridge exchanges a Client~ID and Client
        Secret for a bearer token, then calls the API with
        \code{Authorization: Bearer <token>}, refreshing the token on expiry.
  \item[Fallback — HTTP Basic] Account user/password, used only if client credentials are
        not configured.
\end{description}

\begin{note}
This decision — client credentials primary, Basic fallback, TLS mandatory — is recorded
in \code{docs/adr/0009-smartme-auth-client-credentials.md}. It supersedes the earlier
\code{Authorization: ApiKey} assumption.
\end{note}

\subsection{Configuring it}
Create a \emph{Device (Client Credentials)} OAuth app in the smart-me portal
(\href{https://api.smart-me.com/swagger}{api.smart-me.com}), copy its Client~ID and
secret, and set them in \code{.env}:

\begin{lstlisting}[language=bash]
# smart-me credentials -- LOCAL ONLY, never commit (.env is gitignored, perms 0600)
SMARTME_CLIENT_ID=your-client-id
SMARTME_CLIENT_SECRET=your-client-secret

# Optional HTTP Basic fallback:
# SMARTME_BASIC_USER=your-account-user
# SMARTME_BASIC_PASSWORD=your-account-password
\end{lstlisting}

\begin{warning}
If the client secret is lost or leaked, smart-me does not allow rotating it in place: the
OAuth app must be deleted and recreated. All smart-me traffic is over TLS and hard-fails
if TLS is unavailable.
\end{warning}

\section{MQTT broker connection}
\stub{external broker host/port, optional TLS/auth, broker credentials in \code{.env}.}

\section{Meter-to-topic mapping}
\stub{default scheme, editing, first-run mapping confirmation.}

\section{Poll interval, logging, and retention}
\stub{configurable poll interval, log level, daily rotation and retention window.}
